\documentclass[10pt,a4paper]{amsart}
\usepackage{amsmath,amssymb,amsthm,graphicx}
\usepackage{verbatim}
\newenvironment{reference}{\comment}{\endcomment}
%\newenvironment{reference}{}{}
\newenvironment{slogan}{\comment}{\endcomment}
\newenvironment{history}{\comment}{\endcomment}

% For commutative diagrams we use Xy-pic
\usepackage[all]{xy}

% We use 2cell for 2-commutative diagrams.
\xyoption{2cell}
% Theorem environments.
%
\theoremstyle{plain}
\newtheorem{theorem}[subsection]{Theorem}
\newtheorem{proposition}[subsection]{Proposition}
\newtheorem{lemma}[subsection]{Lemma}

\theoremstyle{definition}
\newtheorem{definition}[subsection]{Definition}
\newtheorem{example}[subsection]{Example}
\newtheorem{exercise}[subsection]{Exercise}
\newtheorem{situation}[subsection]{Situation}

\theoremstyle{remark}
\newtheorem{remark}[subsection]{Remark}
\newtheorem{remarks}[subsection]{Remarks}

\numberwithin{equation}{subsection}

% Macros
%
\def\lim{\mathop{\mathrm{lim}}\nolimits}
\def\colim{\mathop{\mathrm{colim}}\nolimits}
\def\Spec{\mathop{\mathrm{Spec}}}
\def\Hom{\mathop{\mathrm{Hom}}\nolimits}
\def\Ext{\mathop{\mathrm{Ext}}\nolimits}
\def\SheafHom{\mathop{\mathcal{H}\!\mathit{om}}\nolimits}
\def\SheafExt{\mathop{\mathcal{E}\!\mathit{xt}}\nolimits}
\def\Sch{\mathit{Sch}}
\def\Mor{\mathop{\mathrm{Mor}}\nolimits}
\def\Ob{\mathop{\mathrm{Ob}}\nolimits}
\def\Sh{\mathop{\mathit{Sh}}\nolimits}
\def\NL{\mathop{N\!L}\nolimits}
\def\CH{\mathop{\mathrm{CH}}\nolimits}
\def\proetale{{pro\text{-}\acute{e}tale}}
\def\etale{{\acute{e}tale}}
\def\QCoh{\mathit{QCoh}}
\def\Ker{\mathop{\mathrm{Ker}}}
\def\Im{\mathop{\mathrm{Im}}}
\def\Coker{\mathop{\mathrm{Coker}}}
\def\Coim{\mathop{\mathrm{Coim}}}

% Boxtimes
%
\DeclareMathSymbol{\boxtimes}{\mathbin}{AMSa}{"02}

%
% Macros for moduli stacks/spaces
%
\def\QCohstack{\mathcal{QC}\!\mathit{oh}}
\def\Cohstack{\mathcal{C}\!\mathit{oh}}
\def\Spacesstack{\mathcal{S}\!\mathit{paces}}
\def\Quotfunctor{\mathrm{Quot}}
\def\Hilbfunctor{\mathrm{Hilb}}
\def\Curvesstack{\mathcal{C}\!\mathit{urves}}
\def\Polarizedstack{\mathcal{P}\!\mathit{olarized}}
\def\Complexesstack{\mathcal{C}\!\mathit{omplexes}}
% \Pic is the operator that assigns to X its picard group, usage \Pic(X)
% \Picardstack_{X/B} denotes the Picard stack of X over B
% \Picardfunctor_{X/B} denotes the Picard functor of X over B
\def\Pic{\mathop{\mathrm{Pic}}\nolimits}
\def\Picardstack{\mathcal{P}\!\mathit{ic}}
\def\Picardfunctor{\mathrm{Pic}}
\def\Deformationcategory{\mathcal{D}\!\mathit{ef}}
\usepackage{cleveref}
\crefname{originallemma}{original lemma}{original lemmas}
\Crefname{originallemma}{Original lemma}{Original lemmas}
\makeatletter
\newlabel{lemma-monomorphism-separated}{{042N}{source}{}{sourcepointer-1}{}}
\newlabel{lemma-monomorphism-separated@cref}{{[originallemma][0][]042N}{[1][0][]source}{}{}{}}
\makeatother
\crefname{originallemma}{original lemma}{original lemmas}
\Crefname{originallemma}{Original lemma}{Original lemmas}
\makeatletter
\newlabel{lemma-quasi-finite-local}{{040Z}{source}{}{sourcepointer-2}{}}
\newlabel{lemma-quasi-finite-local@cref}{{[originallemma][0][]040Z}{[1][0][]source}{}{}{}}
\makeatother
\crefname{originallemma}{original lemma}{original lemmas}
\Crefname{originallemma}{Original lemma}{Original lemmas}
\makeatletter
\newlabel{morphisms-lemma-etale-over-field}{{02GL}{source}{}{sourcepointer-3}{}}
\newlabel{morphisms-lemma-etale-over-field@cref}{{[originallemma][0][]02GL}{[1][0][]source}{}{}{}}
\makeatother
\crefname{originallemma}{original lemma}{original lemmas}
\Crefname{originallemma}{Original lemma}{Original lemmas}
\makeatletter
\newlabel{morphisms-lemma-quasi-finite-at-point-characterize}{{01TH}{source}{}{sourcepointer-4}{}}
\newlabel{morphisms-lemma-quasi-finite-at-point-characterize@cref}{{[originallemma][0][]01TH}{[1][0][]source}{}{}{}}
\makeatother
\title{A locally finite type monomorphism of algebraic spaces is separated and locally quasi-finite}
\author{The Stacks Project Authors}
\date{stacks-project:spaces-morphisms.tex:local-snapshot-20261012; GNU FDL 1.2 or later}
\begin{document}
\maketitle
\setcounter{subsection}{1}
\noindent\textit{Editorial note.} All delivered mathematical text is the author's own, copied from the pinned Stacks Project chapter file spaces-morphisms.tex; no statement or proof step was rewritten, reordered or removed.\par\medskip
\noindent\textit{Editorial note.} The proof reduces to the affine base and then works with the identity U x\_X U = U x\_Y U coming from the monomorphism hypothesis; that identity is what makes the fibres finite.\par\medskip
\noindent\textit{Editorial note.} The four references of the counted spans are recorded as non-load-bearing navigation pointers; the argument that consumes them stays inside the counted proof span.\par\medskip
\noindent\textit{Editorial note.} Preamble: the three native excerpt spans of sources/stacks-spaces-morphisms/preamble.tex (lines 8-13, 14-18 and 154-217). The chapter file itself carries no preamble text.\par\medskip
\noindent\textit{Allowed standard prerequisites (editorial).}\begin{itemize}
\item A monomorphism of algebraic spaces is separated. Lemma lemma-monomorphism-separated, used at proof lines 5586-5587. Gives the separatedness half of the conclusion.
\item Local quasi-finiteness is local on the base. Lemma lemma-quasi-finite-local, used at proof lines 5588-5591. Reduces to the case of an affine base Y.
\item Fibres of etale morphisms of schemes over a field are finite discrete. Morphisms, Lemma morphisms-lemma-etale-over-field, used at proof lines 5600-5603. Shows the fibres of U -> Y are finite.
\item The criterion for quasi-finiteness at a point via finite fibres. Morphisms, Lemma morphisms-lemma-quasi-finite-at-point-characterize, used at proof lines 5604-5605. Converts finiteness of the fibres into quasi-finiteness of U -> Y.
\end{itemize}\medskip
\ifdefined\hypertarget\hypertarget{sourcepointer-1}{}\fi
\noindent\textit{Source reference.} Original lemma 042N: focused-seeds/AGP/sources/stacks-snapshot-20261011/tags.txt, source lines 5334--5334. Navigation pointer into the Stacks Project; the standard fact it supplies is applied inside the counted proof span and the argument resting on it stays there.\par\medskip
\ifdefined\hypertarget\hypertarget{sourcepointer-2}{}\fi
\noindent\textit{Source reference.} Original lemma 040Z: focused-seeds/AGP/sources/stacks-snapshot-20261011/tags.txt, source lines 5274--5274. Navigation pointer into the Stacks Project; the standard fact it supplies is applied inside the counted proof span and the argument resting on it stays there.\par\medskip
\ifdefined\hypertarget\hypertarget{sourcepointer-3}{}\fi
\noindent\textit{Source reference.} Original lemma 02GL: focused-seeds/AGP/sources/stacks-snapshot-20261011/tags.txt, source lines 3232--3232. Navigation pointer into the Stacks Project; the standard fact it supplies is applied inside the counted proof span and the argument resting on it stays there.\par\medskip
\ifdefined\hypertarget\hypertarget{sourcepointer-4}{}\fi
\noindent\textit{Source reference.} Original lemma 01TH: focused-seeds/AGP/sources/stacks-snapshot-20261011/tags.txt, source lines 2374--2374. Navigation pointer into the Stacks Project; the standard fact it supplies is applied inside the counted proof span and the argument resting on it stays there.\par\medskip
\begin{lemma}
\label{lemma-monomorphism-loc-finite-type-loc-quasi-finite}
Let $S$ be a scheme.
Let $f : X \to Y$ be a morphism of algebraic spaces over $S$.
If $f$ is locally of finite type and a monomorphism, then $f$
is separated and locally quasi-finite.
\end{lemma}

\begin{proof}
A monomorphism is separated, see
Lemma \ref{lemma-monomorphism-separated}.
By
Lemma \ref{lemma-quasi-finite-local}
it suffices to prove the lemma after performing a base change
by $Z \to Y$ with $Z$ affine. Hence we may assume that $Y$ is an
affine scheme. Choose an affine scheme $U$ and an \'etale morphism
$U \to X$. Since $X \to Y$ is locally of finite type the morphism
of affine schemes $U \to Y$ is of finite type.
Since $X \to Y$ is a monomorphism we have $U \times_X U = U \times_Y U$.
In particular the maps $U \times_Y U \to U$ are \'etale.
Let $y \in Y$. Then either $U_y$ is empty, or
$\Spec(\kappa(u)) \times_{\Spec(\kappa(y))} U_y$
is isomorphic to the fibre of $U \times_Y U \to U$ over $u$ for
some $u \in U$ lying over $y$. This implies that the fibres of
$U \to Y$ are finite discrete sets (as $U \times_Y U \to U$
is an \'etale morphism of affine schemes, see
Morphisms, Lemma \ref{morphisms-lemma-etale-over-field}).
Hence $U \to Y$ is quasi-finite, see
Morphisms, Lemma \ref{morphisms-lemma-quasi-finite-at-point-characterize}.
As $U \to X$ was an arbitrary \'etale morphism with $U$ affine
this implies that $X \to Y$ is locally quasi-finite.
\end{proof}

\end{document}
